\documentclass[10pt]{extarticle}
\usepackage[utf8]{inputenc}
\usepackage[brazil]{babel}
\usepackage{graphicx}
\usepackage{amsmath}
\usepackage{amssymb}
\usepackage{xcolor}
\usepackage{listings}
\usepackage{paralist}
\usepackage{algorithmic}
\usepackage{xurl}
\usepackage{titlesec}
\usepackage[left=1.5cm, right=1.5cm, top=1.5cm, bottom=2cm, twocolumn]{geometry}
\usepackage{hyperref}
\usepackage{caption}
\usepackage{titling}
\usepackage{fancyhdr}
\usepackage{float}
\usepackage{booktabs}

\usepackage{abstract}

\usepackage{tikz}
\usetikzlibrary{automata,arrows,calc,shapes,decorations.pathreplacing}
\tikzset{->,>=stealth',auto,node distance=2cm,thick,initial text=}
\tikzset{
    position/.style args={#1:#2 from #3}{
        at=(#3.#1), anchor=#1+180, shift=(#1:#2)
    }
}

\renewcommand\lstlistingname{Código}
\newcommand{\lstnumbername}{Linha}
\newcommand{\lstlistingautorefname}{Código}
\lstdefinestyle{one_column}{
    backgroundcolor=\color{yellow!10!white},
    commentstyle=\color{cyan!50!black},
    keywordstyle=\color{magenta},
    numberstyle=\tiny\color{gray},
    stringstyle=\color{violet},
    basicstyle=\ttfamily\scriptsize,
    breakatwhitespace=false,
    breaklines=true,
    captionpos=b,
    keepspaces=true,
    numbers=left,
    numbersep=5pt,
    showspaces=false,
    showstringspaces=false,
    showtabs=false,
    tabsize=1,
    frame=single,
    language=C++,
    xleftmargin=15pt,
    linewidth=255pt,
    escapechar=@,
}
\lstset{style=one_column}

\pretitle{%
    \vspace{-1.5cm}
    \centering
    \color{cyan!50!black}
    \huge\bfseries
}%
\posttitle{%

    \color{cyan!50!black}
    \vspace{-0.4cm}
    \noindent\vrule height 2.5pt width 0.49\textwidth
    \vskip 0cm
}

\titleformat{\section}
{\color{cyan!50!black}\normalfont\Large\bfseries}
{\color{cyan!50!black}\thesection}{1em}{}

\titleformat{\subsection}
{\color{cyan!50!black}\normalfont\Large\bfseries}
{\color{cyan!50!black}\thesection}{1em}{}

\renewcommand{\abstractname}{Resumo}

\renewcommand{\abstractnamefont}{%
    \bfseries
    \color{cyan!50!black}
}

% TODO: confirm final title, authors and affiliations.
\title{Filas de Prioridade Monótonas para a Otimização de Dijkstra}
\author{\color{cyan!50!black}\bfseries Lucas Pavanelli % TODO: instituição / coautores
}
\date{\vspace{-1.0cm}}

\begin{document}

\titlespacing*{\section}
{0pt}{3pt}{4pt}
\titlespacing*{\subsection}
{0pt}{3pt}{4pt}

\maketitle

\begin{abstract}
% TODO: escrever o resumo final (problema + algoritmos + principal achado experimental).
Este artigo compara variantes de filas de prioridade monótonas para o
algoritmo de Dijkstra em grafos com pesos inteiros não-negativos: o
\textit{binary heap} clássico, o \textit{radix heap} de um nível e o
\textit{radix heap} de dois níveis. Discutimos as complexidades
$\mathcal{O}((m+n)\log n)$, $\mathcal{O}(m + n\log C)$ e
$\mathcal{O}(m + n\log C / \log\log C)$, descrevemos as implementações em C++
e avaliamos experimentalmente onde cada estrutura é vantajosa, com foco em
aplicações de programação competitiva.
\end{abstract}

% Corpo do artigo. Cada seção abaixo é um esqueleto a ser preenchido nos
% blocos de escrita (ver docs/timeline.md). Mantenha o estilo denso de duas
% colunas do template (~3000 palavras em 2 páginas).

\section{Introdução}
% TODO (bloco Qua Jun 24): caminhos mínimos em programação competitiva, por que
% a escolha da fila importa, visão geral do artigo de referência (Costa et al.).
O problema do caminho mínimo de fonte única (\textit{single-source shortest
path}, SSSP) é recorrente em programação competitiva. Em grafos com pesos
inteiros não-negativos, o algoritmo de Dijkstra resolve o problema, e seu
desempenho depende fortemente da fila de prioridade utilizada. Neste artigo
revisitamos filas de prioridade \textit{monótonas} --- aquelas em que as chaves
extraídas nunca decrescem --- seguindo o levantamento de Costa et al.
\cite{costa2025monotone}.

\section{Algoritmos}
% TODO (bloco Qua Jun 24): binary heap (baseline), radix heap de um nível e de
% dois níveis; explicar o truque XOR/dígito base-K.
Comparamos três filas de prioridade. O \textit{binary heap} serve de linha de
base. O \textit{radix heap} de um nível distribui as chaves em \textit{buckets}
indexados pelo bit mais significativo em que diferem do mínimo corrente. O
\textit{radix heap} de dois níveis generaliza essa ideia para dígitos na base
$K$, reduzindo o número de níveis que um elemento percorre.

\begin{table}[H]
    \centering
    \caption{Complexidade de Dijkstra por fila de prioridade. $C$ é o maior
    peso de aresta.}
    \label{tab:complexidade}
    \footnotesize
    \begin{tabular}{@{}ll@{}}
        \toprule
        Fila de prioridade & Complexidade \\
        \midrule
        Binary heap & $\mathcal{O}((m+n)\log n)$ \\
        Radix heap (1 nível) & $\mathcal{O}(m + n\log C)$ \\
        Radix heap (2 níveis) & $\mathcal{O}(m + n\log C / \log\log C)$ \\
        Radix + Fibonacci heap & $\mathcal{O}(m + n\sqrt{\log C})$ \\
        \bottomrule
    \end{tabular}
\end{table}

% A variante com Fibonacci heap é apenas mencionada (teórica); não a implementamos.

% Foreshadowing (expandir no bloco Qua Jun 24): a Tabela~\ref{tab:complexidade}
% troca o fator $\log n$ por operação do binary heap por $\log C$. Logo, espera-se
% que os radix heaps vençam quando $\log C < \log n$ (pesos pequenos e/ou grafos
% grandes) e percam quando $C$ é muito grande --- hipótese que os experimentos
% testam.

\section{Implementação}
% TODO (bloco Qui Jun 25): trechos de código curtos -- estrutura de buckets do
% radix heap e o índice de bucket via XOR (one-level) / dígito base-K (two-level).
% Exemplo de inclusão de código:
% \begin{lstlisting}[caption={Índice de bucket no radix heap de dois níveis.}]
% ... colar trecho de include/radix_heap.hpp ...
% \end{lstlisting}

\section{Experimentos}
% TODO (bloco Qui Jun 25): descrever metodologia e famílias.
% Metodologia: 5 fontes x 3 repetições por configuração, mediana reportada,
% semente fixa (seed=1); experiments/full_sweep.sh cobre as 4 famílias x
% C in {10, 1e3, 1e6, 1e9}. Hardware/compilador em manifest.txt.
% Famílias e por que cada uma (variam densidade, estrutura e dispersão de
% distâncias): grid (planar, esparso, distâncias muito dispersas), sparse
% (aleatório, m ~ 3n), dense (Erdős–Rényi, m ~ n^2, portanto n pequeno),
% layered (DAG em camadas, distâncias concentradas por camada).
%
% Tabela de runtime (Tabela~\ref{tab:resultados}) e contadores de operação
% (Tabela~\ref{tab:opcounts}) são geradas por scripts/analyze_results.py a
% partir do results.csv mais recente; regenere após cada full_sweep.
% Generated by scripts/analyze_results.py -- do not edit by hand.
\begin{table}[H]
    \centering
    \caption{Tempo mediano de execução (ms) por fila de prioridade. Melhor valor por linha em \textbf{negrito}.}
    \label{tab:resultados}
    \footnotesize
    \begin{tabular}{@{}llrrr@{}}
        \toprule
        Família & $C$ & binary & radix-1 & radix-2 \\
        \midrule
        dense ($n${=}2k) & 10 & 1.03 & \textbf{0.93} & 0.96 \\
         & 1\,000 & 1.30 & 1.24 & \textbf{1.08} \\
         & 1\,000\,000 & \textbf{1.35} & 1.98 & 1.73 \\
         & 1\,000\,000\,000 & \textbf{1.36} & 2.01 & 1.92 \\
        \midrule
        grid ($n${=}90k) & 10 & 3.73 & 3.10 & \textbf{2.78} \\
         & 1\,000 & \textbf{4.38} & 10.88 & 5.18 \\
         & 1\,000\,000 & \textbf{4.51} & 12.94 & 12.01 \\
         & 1\,000\,000\,000 & \textbf{4.43} & 13.18 & 12.49 \\
        \midrule
        grid ($n${=}360k) & 10 & 17.52 & 14.73 & \textbf{13.28} \\
         & 1\,000 & 21.31 & 41.48 & \textbf{21.22} \\
         & 1\,000\,000 & \textbf{21.43} & 59.59 & 51.41 \\
         & 1\,000\,000\,000 & \textbf{21.57} & 57.51 & 55.27 \\
        \midrule
        layered ($n${=}120k) & 10 & 4.41 & 2.99 & \textbf{2.82} \\
         & 1\,000 & 6.02 & 8.26 & \textbf{4.74} \\
         & 1\,000\,000 & \textbf{5.66} & 13.67 & 12.54 \\
         & 1\,000\,000\,000 & \textbf{6.01} & 14.16 & 13.79 \\
        \midrule
        sparse ($n${=}250k) & 10 & 19.47 & \textbf{11.47} & 11.76 \\
         & 1\,000 & 25.35 & 15.19 & \textbf{13.89} \\
         & 1\,000\,000 & \textbf{26.03} & 44.13 & 31.29 \\
         & 1\,000\,000\,000 & \textbf{27.15} & 45.17 & 41.31 \\
        \bottomrule
    \end{tabular}
\end{table}


% Generated by scripts/analyze_results.py -- do not edit by hand.
\begin{table}[H]
    \centering
    \caption{Contadores de operação por família (média sobre as execuções; idênticos entre as filas). push/$n$ = relaxações bem-sucedidas por vértice; \textit{stale}\,\% = fração de remoções obsoletas (\textit{lazy deletion}).}
    \label{tab:opcounts}
    \footnotesize
    \begin{tabular}{@{}lrrrr@{}}
        \toprule
        Família & $n$ & $m$ & push/$n$ & \textit{stale}\,\% \\
        \midrule
        dense & 2\,000 & 1\,202\,346 & 5.06 & 80.2 \\
        grid & 90\,000 & 358\,800 & 1.31 & 23.5 \\
        grid & 360\,000 & 1\,437\,600 & 1.31 & 23.6 \\
        layered & 120\,001 & 598\,800 & 0.98 & 37.2 \\
        sparse & 250\,000 & 750\,000 & 1.07 & 15.8 \\
        \bottomrule
    \end{tabular}
\end{table}


% A Tabela~\ref{tab:opcounts} expõe as causas estruturais usadas na Análise:
% push/$n$ (atividade de decrease-key) e \textit{stale}\% (remoções obsoletas
% sob lazy deletion). Note o $n$ pequeno de dense (m ~ n^2 limita n) e a baixa
% fração stale de sparse.

\section{Análise}
% NOTAS PARA ESCRITA (bloco Sex Jun 26). Expandir em prosa densa.
% (Atenção: radix vence com C PEQUENO; perde com C GRANDE.)
%
% Resumo do achado: a C=10 todas as famílias favorecem um radix (em geral o de
% dois níveis); o cruzamento ocorre entre C=1e3 e C=1e6; a partir daí o binary
% heap vence. Ordem da vantagem do radix a C=10:
% sparse 1.70x > layered 1.57x > grid 1.32x > dense 1.11x.
%
% TRÊS ALAVANCAS que explicam o resultado:
%   1. n (tamanho do heap): custo do binary ~ (pushes+pops)*log n. n maior =>
%      mais a economizar com push O(1) do radix. (ver Tabela~\ref{tab:opcounts})
%   2. stale% (lazy deletion): muitas remoções obsoletas inflam o fator constante
%      do radix (churn nos buckets).
%   3. dispersão de distâncias: define o quão forte o termo n*log C pesa; quanto
%      mais larga/uniforme a faixa de distâncias, mais cedo o radix degrada com C.
%
% POR FAMÍLIA:
%   - dense: vence menos (1.11x). m ~ n^2 limita n a 2000 => log n ~ 11, barato
%     para o binary; além disso a maior fração stale (Tabela~\ref{tab:opcounts}).
%   - grid: vitória moderada a C pequeno, mas DEGRADA MAIS RÁPIDO com C ---
%     distâncias muito dispersas/uniformes fazem o n*log C explodir; o radix de
%     1 nível colapsa a C=1e3 (grid n=360k: 41ms vs 21ms do de 2 níveis).
%   - layered: vitória forte (1.57x) e robusta a C --- distâncias concentradas
%     por camada (frente de onda monótona) minimizam redistribuição, mesmo com
%     menos pushes (DAG).
%   - sparse: vence mais (1.70--1.83x) --- maior n (log n ~ 18) e menor fração
%     stale; toda alavanca aponta para o radix.
%
% DOIS NÍVEIS vs UM NÍVEL: o radix de dois níveis domina o de um nível em quase
% toda a tabela; o caso ilustrativo é grid C=1e3 (colapso do de 1 nível). O
% fator log log C amortece o crescimento com C.

\section{Conclusão}
% TODO (bloco Sex Jun 26): conclusões práticas para programadores competitivos.

% Trabalhos futuros: avaliar as filas em redes reais (instâncias DIMACS de
% estradas dos EUA), cujos pesos correspondem a distâncias/tempos --- $C$
% moderado em grafos esparsos e quase planares, próximo do caso favorável aos
% radix heaps. Exigiria um leitor de arquivos \texttt{.gr}/\texttt{.co}. Também
% fica em aberto a variante com Fibonacci heap ($\mathcal{O}(m + n\sqrt{\log C})$).


\bibliographystyle{plain}
\bibliography{references}

\end{document}
